\begin{theorem}[The pair entries along a rank-one supercharge]
\label{thm:c11-twistentries}
Let $\xi=u\otimes z\in S_+\otimes W$ be of rank one, with $u'$, the twisted letters and their representatives of Remark \ref{rem:c11-twist}, and let $\psi=0$. Normalise the letters as the components along the weight vectors of a Cartan subalgebra of $\mathfrak{so}(V)\oplus\mathfrak{so}(W_{\mathbb R}^\perp)$ for which $u$, $u'$ and $z$ are weight vectors, read the outputs on the dual weight vectors, take $b$ as the component $F(d_1,d_2)$ of the field strength on the two null vectors $d_1$, $d_2$ of $V\otimes\mathbb C$ along which $\partial_{\dot\alpha}$ differentiates, and expand the output momenta in the null coordinates; this fixes each entry up to one nonzero factor for each letter and each output. Reduce modulo $\operatorname{Im}Q^\xi_0$, total derivatives kept: the derivatives $u^\alpha\sigma^m_{\alpha\dot\alpha}\partial_m$ are $Q^\xi_0$-exact on the letters, and so is, up to terms with antifields, the holomorphic Dirac component of a derivative on an output $\partial_{\dot\alpha}c$. Then the arity-one pair entries of $Q_1$ are those of Table \ref{tab:c11-twistentries}, and:
\begin{itemize}
\item (i) the entry $\beta\gamma$, the kernel of $(\lambda|_{S_+\otimes W},A|_{W_{\mathbb R}^\perp})$ on the leg set $(\lambda|_{S_-\otimes W^*})^2$, vanishes for every $\xi\in S_+\otimes W$, the placements on a line and at the special edge cancelling chain by chain; the one-body terms and $L_1^{(0)}$ have outputs of other Weyl weights or with a spectator letter, so the Konishi entry of Remark \ref{rem:c11-twist} is $0$ at $\psi=0$; the entry $bc$ with $b$ represented by $(\xi,\lambda^*)$ and the entry $\gamma\gamma$ vanish as well;
\item (ii) the entries $\beta\beta$ and $\beta b\to\gamma\gamma$ carry the colour and derivative structure of [BK, (3.67)] and of the last term of [BK, (3.69)]: the output $\gamma^K$, resp. the outputs $\gamma^J$, $\gamma^K$, with the indices complementary to those of the letters $\beta$, the same value in the two colour orders, and, for $\beta b\to\gamma\gamma$, the antisymmetric combination of the two holomorphic derivatives;
\item (iii) the entry $b\gamma$ has different values in the two colour orders, each with a term in $\log\tfrac32$, and the entry $\beta b\to\beta\,\partial c$ is nonzero in one colour order;
\item (iv) every derivative in the entries is holomorphic.
\end{itemize}
Statements (i)--(iv) are independent of the normalisation factors; the values of the individual entries depend on the scheme (Proposition \ref{prop:c11-freedom} (iv), Remark \ref{rem:c11-twist}).
\end{theorem}

\begin{tabular}{l|l|l|l|l}
entry & pair $\to$ outputs & order $(Y_1,Y_2)$ & order $(Y_2,Y_1)$ & Dirac components\\
\hline
$\beta\gamma$ [BK, (3.66)] & $(\beta_I,\gamma^J)\to\partial c\,\partial c$ & $0$ & $0$ & $0$, $0$\\
$bc$ [BK, (3.64)], $b\leftrightarrow(\xi,\lambda^*)$ & $(c,b)\to\partial c\,\partial c$ & $0$ & $0$ & $0$, $0$\\
$\beta\beta$ [BK, (3.67)] & $(\beta_I,\beta_J)\to\gamma^K\,\partial c$ & $\tfrac23$ & $\tfrac23$ & $\tfrac13-2L$, $\tfrac13-2L$\\
$b\gamma$ [BK, (3.68)] & $(b,\gamma^I)\to\gamma^I\,\partial c$ & $\tfrac{64}3-128L$ & $32+64L$ & $-\tfrac{64}3-128L$, $-\tfrac{32}3+64L$\\
$\beta b\to\beta\,\partial c$ [BK, (3.69)] & $(\beta_I,b)\to\beta_I\,\partial c$ & $-\tfrac{64}3$ & $0$ & $0$, $0$\\
$\beta b\to\gamma\gamma$ [BK, (3.69)] & $(\beta_I,b)\to\gamma^J\gamma^K$ & $\tfrac83\big(\partial_2\gamma^J\partial_1\gamma^K-\partial_1\gamma^J\partial_2\gamma^K\big)$ & the same & \\
$\gamma\gamma$ [BK, (3.77)] & $(\gamma^I,\gamma^J)\to$ twisted letters & $0$ & $0$ & \\
\end{tabular}
\tabcaption{The arity-one pair entries of $Q_1$ along $\xi=u\otimes z$ at $\psi=0$, reduced modulo $\operatorname{Im}Q^\xi_0$, in the normalisation of Theorem \ref{thm:c11-twistentries}: coefficients of $\frac{\hbar}{16\pi^2}$ times the colour factor $(1\otimes\mathrm{ad}_{Y_2}\mathrm{ad}_{Y_1})(-,-)^{-1}$ of the colour order $(Y_1,Y_2)$, $Y_1$ the first output listed, with $L=\log\tfrac32$. The third and fourth columns give the coefficient of the holomorphic derivative on the output other than $\partial c$, and for $\beta b\to\gamma\gamma$ the derivative structure, $\partial_1$ and $\partial_2$ being the derivatives along $d_1$ and $d_2$; the last column gives, in the two orders, the coefficient of the holomorphic Dirac component on the output $\partial c$, which is $Q^\xi_0$-exact up to terms with antifields; $K$ is complementary to $I,J$, and $J,K$ to $I$.\label{tab:c11-twistentries}}

\begin{proof}
The entries are the kernels on the leg sets $(\lambda|_{S_-\otimes W^*})^2$ of $(\lambda|_{S_+\otimes W},A|_{W_{\mathbb R}^\perp})$; $A|_{W_{\mathbb R}^\perp}\lambda|_{S_-\otimes W^*}$ of $(\lambda|_{S_+\otimes W},\lambda|_{S_+\otimes W})$ and of $(A|_V,A|_{W_{\mathbb R}^\perp})$ with one derivative on $A|_V$; $\lambda|_{S_+\otimes W}\lambda|_{S_-\otimes W^*}$ and $(A|_{W_{\mathbb R}^\perp})^2$ of $(\lambda|_{S_+\otimes W},A|_V)$ with one derivative on $A|_V$; and $(\lambda|_{S_-\otimes W^*})^2$ of $(c,\lambda^*|_{S_+\otimes W^*})$ (Tables \ref{tab:c11-defectpairs} and \ref{tab:c11-defectcounts}); all have superficial degree $0$, so they carry no poles and no finite special vertex (Proposition \ref{prop:c11-oddbound}), and Theorem \ref{thm:c11-oddsimplex} evaluates them exactly, with the letter vectors contracted exactly in the weight basis (Remark \ref{rem:c11-oddmethod}). The colour factors are those of Lemma \ref{lem:c10-colour} with the output of a term $T_5$ read as an output of a cubic vertex. (i) The vanishing of $\beta\gamma$ for general $\xi$ and of $bc$ holds in both evaluations. The outputs $(\lambda|_{S_-\otimes W^*})^2$ have Weyl weight $3$, while $L_1^{(0)}$ preserves the Weyl weight $\tfrac52$ of the Konishi descendant, a one-body term keeps the other letter as a spectator, and the arity-two kernels have outputs of degree $-1$; hence the coefficient of $\int\omega\,(\bar\lambda_{\dot\alpha i}z^i,\bar\lambda^{\dot\alpha}{}_jz^j)$ in $\hbar Q^\xi_1$ of the Konishi descendant is the entry $\beta\gamma$. $\gamma\gamma$ is Proposition \ref{prop:c11-plane}. (ii)--(iv) are read off the table. Rescaling a weight vector rescales its letter by the inverse factor and its output by the factor, which multiplies an entry by one product of factors, the same for the two colour orders, for the reduced and the Dirac coefficients, and for the rational part and the part in $\log\tfrac32$; the vanishing, the equality of the two orders, the indices and the derivative structure are therefore independent of the normalisation.
\end{proof}

\begin{theorem}[What the comparison with the holomorphic twist tests]
\label{thm:c11-twistcheck}
Let $\xi=u\otimes z$ be of rank one, let $\Phi$ be the quasi-isomorphism from the cohomology of $Q^\xi_0$ on local operators to that of the free holomorphic differential on the letters $c$, $\gamma$, $\beta$, $b$, $\partial_{\dot\alpha}$ of [BK], transported through [ESW, Thm 10.9] (remark after Theorem \ref{thm:c8-twist}), and let $D$ be the difference, read through $\Phi$, of the one-loop quantisation $r_\xi(\hbar F_1)$ of the twisted theory restricted from the equivariant one and the holomorphic quantisation of [BK, (2.42), footnote 22] and [W18, Thm 3.4] (B. R. Williams, Renormalization for holomorphic field theories, arXiv:1809.02661). If $D\in\mathbb C\,S^\xi+\operatorname{Im}Q^\xi_0$, then $\Phi\circ[Q^\xi_1]\circ\Phi^{-1}$ is the class of the one-loop differential of [BK, (3.7), (3.64)--(3.69)], and:
\begin{itemize}
\item (a) the one-body terms of $Q^\xi_1$ on the twisted letters are $a\,Q^\xi_0+[Q^\xi_0,N]$ for a number $a$ and a derivation $N$ of the letters, zero on cohomology;
\item (b) on the pairs from $\{c,\gamma\}$ and on $(\beta,c)$ the two-body kernels lie, on the twisted letters, in $\operatorname{Im}Q^\xi_0$ plus terms cancelled on cohomology by the one-body terms, as in [BK, (3.11)]; on $(\gamma,\gamma)$ they vanish (Proposition \ref{prop:c11-plane});
\item (c) on the pairs $\beta\gamma$, $\beta\beta$, $b\gamma$, $\beta b$, $bb$ and $bc$ the kernels, reduced modulo $\operatorname{Im}Q^\xi_0$ as in Theorem \ref{thm:c11-twistentries}, have the colour factors $(1\otimes\mathrm{ad}_{Y_2}\mathrm{ad}_{Y_1})(-,-)^{-1}$, which are $\tfrac12f_{ACD}f_{BCE}$ of [BK] up to normalisation, and holomorphic derivatives $\partial_{\dot\alpha}=u'^\alpha\sigma^m_{\alpha\dot\alpha}\partial_m$ on the outputs;
\item (d) their coefficients satisfy, modulo the finite parts of Proposition \ref{prop:c11-freedom} (iv), the relations [BK, (3.59)--(3.63)], which are the order-$\hbar$ part of Corollary \ref{cor:c8-homotopy} (ii) at $(\xi,R)$ for the elements $R$ of the tree-level commutant $\mathfrak{psl}(3|3)$ of $\xi$, restricted by $r_\xi$;
\item (e) every statement of the holomorphic twist about $[Q_1]$ on classes of several letters is the image under $\Phi$ of $[Q^\xi_1]$, computed by the pair rule from the kernels of (c), with the normalisation of $\hbar$ relative to the unit fixed by $\Phi$.
\end{itemize}
For a general $D$ the two class maps differ by $[\{D,-\}]$, and (a)--(e) compute that class.
\end{theorem}

\begin{proof}
Two quantisations of the twisted theory differing by a $Q^\xi_0$-closed $D$ of degree $0$ have class maps differing by $[\{D,-\}]$, which vanishes for $D\in\mathbb C\,S^\xi+\operatorname{Im}Q^\xi_0$, since $\{S^\xi,-\}=Q^\xi_0$ and $\{Q^\xi_0\chi,-\}=\pm[Q^\xi_0,\{\chi,-\}]$ (Theorem \ref{thm:c8-twist} and the remark after it). (a) The one-body terms are brackets with local functionals (Theorem \ref{thm:c11-classes} (iii)), and the condition on $D$ makes them $aQ^\xi_0$ plus commutators with $Q^\xi_0$ on the twisted letters. (b), (c) The pair rule [BK, (3.7)--(3.11)], the colour factors of Lemma \ref{lem:c10-colour} with two outputs, and the exactness of the antiholomorphic derivatives, $\{Q,\bar Q^1_{\dot\alpha}\}=P_{-\dot\alpha}$ [BK, (2.5)]. (d) Corollary \ref{cor:c8-homotopy} (ii) at $(\xi,R)$ with $[\xi,R]=0$, at order $\hbar$, mapped by $r_\xi$ (Lemma \ref{lem:c8-ev}). (e) $[Q^\xi_1]$ is determined by its values on representatives, polynomials in the letters, through the pair rule.
\end{proof}

\begin{remark}[Method of computation]
\label{rem:c11-oddmethod}
Each defect cycle is evaluated exactly as in Remark \ref{rem:c11-method}: the Gaussian loop integral and the integrals of \eqref{eq:c11-oddsimplex} over the ordered subsimplices of the region where $\alpha_\sigma$ is the smallest coordinate, the letter derivatives through Proposition \ref{prop:c11-genfun} with $n+1$ edges, the position along $\eta$ and $K_b$ through the injected momentum $r$, the fibre algebra in an explicit realisation of $V_{10}$ and $\Sigma$ over $\mathbb Q(i)$ with the vertices $S^{(1)}(x)$ of Chapter 2 and the sign of Definition \ref{def:c11-defect} (ii), letter vectors with fractional components contracted exactly, and the colour words by Lemma \ref{lem:c10-colour}; the twisted letters are built from the weights of a Cartan subalgebra. A second evaluation works in momentum space with an independently written realisation of $S^{(1)}(x)$ in the components of Definition \ref{def:c8-twocomp}, the special edge as the factor $e^{-\varepsilon|p|^2}-1$, Feynman parameters for the other edges and numerical integration at several values of $\varepsilon$ with thirty digits, extrapolated to $\varepsilon^0$ with the poles fitted; along $K_b$ it evaluates the cycles at the endpoints $T=1,4,16$. The two evaluations give the same operators $C_\xi$, $Y_\eta$, $C_K$ and the same arrows of Lemma \ref{lem:c11-lines}, the same composite line along $\xi$, the same free directions of Proposition \ref{prop:c11-freexi}, the same kernels $\Upsilon_\xi$ on every leg set with at most two outputs of the pairs of field and ghost letters without derivatives and with one letter derivative, the same coefficients of Theorem \ref{thm:c11-bilinears} for both chiralities, the same entries $\beta\gamma$, $\beta\beta$ and $bc$, and the same translation-invariant kernels along $\eta$ and $K_b$ where both are evaluated; both reproduce Example \ref{ex:c11-cubic}, Proposition \ref{prop:c11-plane} and the value along $P_b$ of Corollary \ref{cor:c11-DPdefect}. The first evaluation gives the values along $\Delta$ of Corollary \ref{cor:c11-DPdefect}, extends $\Upsilon_\xi$ to three letter derivatives without the finite special vertices, and gives the entries of Table \ref{tab:c11-twistentries} and the identity (iv) of Theorem \ref{thm:c11-bilinears}; the second gives the closed forms of Theorem \ref{thm:c11-oddsimplex} (iv) for $n_t=2,\dots,6$, the finite special vertices on the pairs of gaugino letters, the pole identity \eqref{eq:c11-pole} on these pairs and on the pair of Example \ref{ex:c11-cubic} with the vacuum-pole terms, the position identity of Theorem \ref{thm:c11-oddsimplex} (ii), and, along $K_b$, $T\partial_T$ of the defect cycles of the intrinsic vertex $S^{(1)}(K_b)$ on $(A|_V,A|_V)$ with the output $A|_V$, with the vertex on a line and at the special edge, decreasing at least as $\varepsilon$ at the three endpoints, as Theorem \ref{thm:c11-oddsimplex} (iii) states. Every integral is a four-dimensional Schwinger integral with the cutoff $\varepsilon$ on every edge.
\end{remark}

\begin{remark}[The data that complete the one-loop components]
\label{rem:c11-next}
The one-body kernels of $Q_1$, $S_1$ and $K_1$ are given by the action side of the arity-one equations. Along $\xi$ it is \eqref{eq:c11-arity1xi}: the quadratic part of $\hbar\,\mathrm{Ob}^{(1)}(\xi)$ in the sectors $(\lambda,A|_V)$, $(\lambda,A|_{W_{\mathbb R}^\perp})$ and $(c,\lambda^*)$, the bracket $\{S^{(1)}(\xi),\hbar F_1^{(0)}\}$ with the finite counterterm of Theorem \ref{thm:c4-finite}, a particular solution, whose coefficient of $T_2$ gives $t_\perp$ and with it $c_Q$ of Theorem \ref{thm:c11-bilinears}, and the cubic part, which fixes $\theta^\xi_V-\theta^\xi_\perp$; the same computation along $\eta$ and $K_b$ gives the remaining one-body kernels. The next computations of the two-body kernels are the finite special vertices on the leg sets with $D\ge2$, with their vacuum-pole terms; the arity-two kernels along $(\xi,\xi)$, which give the entry $bc$ with $b$ represented by the self-dual field strength; and the entry $bb$ on the pair $(A|_V,A|_V)$ with two derivatives. With the one-body kernels, the relations (R2)--(R6) of Proposition \ref{prop:c11-relations} become finite computations on the scalar bilinears, with the cocycle finite part of Theorem \ref{thm:c11-bilinears} (iii) extended beyond the two-field sector; (R3) on the trace bilinear is the Konishi relation between $Q_1$, $S_1$ and $D_1=-48$ units. The class-level comparison of Theorem \ref{thm:c11-twistcheck} uses the transport map $\Phi$ built from [ESW, Thm 10.9], the condition on $D$ and the relations [BK, (3.59)--(3.63)] applied to the entries of Table \ref{tab:c11-twistentries}. The inner-product relations of Proposition \ref{prop:c11-S} (iv) and Corollary \ref{cor:c11-K} use $\mathrm{Gram}_{\rm 1}$ from the two-point graphs of Definition \ref{def:c8-2pt}, and $K_1\cong-(b,x)D_1$ on the families uses the rotation component $M_1$.
\end{remark}

